%% kapa: exam=1 round=37 year=2026 period=1 form=A subject=real_estate subject_ko=부동산학원론 seq=4 number=84
%% source: Q-Net 감정평가사 기출문제 (cst00309) artlSeq=5258924 file=raw/1008/5258924/2026년 제37회 감정평가사 1차 1교시.pdf
%% conversion: kordoc 4.24.51 -> tools/gen1.py (auto); review=proofread-sonnet; flags=none
\begin{question}{84}
토지의 특성과 내용에 관한 설명으로 옳은 것은 모두 몇 개인가?
\begin{framedtext}
○ 토지의 개별성은 부동산시장을 불완전경쟁시장으로 만드는 요인이다. \par
○ 토지의 용도가 다양하므로 부동산활동에서 최유효이용에 대한 분석이 필요하다. \par
○ 토지의 영속성으로 인해 감정평가에 있어 장래의 동향에 대한 분석이 필요하다. \par
○ 토지는 공간적으로 연결되어 있으므로 부동산활동을 위해 인위적으로 구획하여 등록하는 것이 필요하다. \par
○ 하나의 토지에 소유자와 용익물권자가 각각 있는 경우 각 권리별로 감정평가해야 할 필요가 있다. \par
\end{framedtext}
\begin{choices}
  \item 1개
  \item 2개
  \item 3개
  \item 4개
  \item 5개
\end{choices}
\end{question}
\answer{5}
